\documentclass[10pt,a4paper]{amsart}
\usepackage[margin=19mm]{geometry}
\usepackage{amsmath,amssymb,mathtools}
\usepackage[scr=rsfs,cal=euler]{mathalfa}
\usepackage{xfrac}
\usepackage[unicode,hidelinks,bookmarks=false]{hyperref}
\newcommand{\ie}{i.e.\ }
\newcommand{\cf}{cf.\ }
\newcommand{\resp}{respectively\ }
\newcommand{\Subsection}{\subsection}
\providecommand{\nobreakdash}{\nobreak\hbox{-}}
\setlength{\emergencystretch}{2em}
\multlinegap0pt
\newcommand{\Pfour}{\mathbb{P}^4}
\usepackage{amssymb}
\usepackage{xcolor}
\usepackage[shortlabels]{enumitem}
\usepackage{mathtools}
\usepackage{cancel}
\usepackage{float}
\usepackage{caption}
\newtheorem{lem}{Lemma}
\newcommand{\Aut}{\textup{Aut}}
\newcommand{\GL}{\textup{GL}}
\newcommand{\PGL}{\textup{PGL}}
\newcommand{\Pone}{\mathbb{P}^1}
\newcommand{\Pthree}{\mathbb{P}^3}
\newcommand{\Ptwo}{\mathbb{P}^2}
\newcommand{\VPthree}{V_{\Pthree}}
\newcommand{\VPtwo}{V_{\Ptwo}}
\newcommand{\kk}{\mathbf{k}}
\newcommand{\sub}{\subseteq}
\title{Weak Fano threefolds arising as the blowup of a hyperquadric in \texorpdfstring{$\Pfour$}{P4} along a curve}
\author{Anne Schnattinger}
\date{Source: arXiv:2512.20315v2 (CC BY 4.0)}
\begin{document}
\maketitle

\noindent\emph{Source and licence.} Weak Fano threefolds arising as the blowup of a hyperquadric in \texorpdfstring{$\Pfour$}{P4} along a curve. Version arXiv:2512.20315v2 (CC BY 4.0), distributed by arXiv under the Creative Commons Attribution 4.0 International licence, http://creativecommons.org/licenses/by/4.0/, as recorded in the arXiv licence metadata for this exact version and shown on the abstract page https://arxiv.org/abs/2512.20315v2. Full licence notice: \texttt{licenses/arxiv-2512.20315-CC-BY-4.0.txt}.\par\smallskip

\noindent\textit{Editorial scope.} Counted unit: the article's lem lem:automorphisms\_quadric\_cone (source0.tex lines 1392--1394). The article's complete original proof of it is delivered with it (source0.tex lines 1396--1423, one \texttt{proof} environment of 28 lines). No mathematics is written, repaired or re-derived: the statement and the proof are the article's own lines, in the article's own notation, and no retained line is substituted or re-flowed.  No further source-local context is needed: every cross-reference of every retained line resolves inside the two delivered spans.  Nothing was omitted from any retained span, so the package carries no omission record.  No retained line contains a citation, so the wrapper carries no bibliography and no bibliography entry.  No retained line contains a figure environment, an image-inclusion command or drawing code, so no image is delivered and the package declares no asset dependency.  Editorial formatting only, in addition: an AMS \texttt{amsart} preamble that restates the article's own declarations and macros used by the retained lines, namely the environments \texttt{lem} on separate counters, together with the standard packages those lines need.  The retained environments print in this document as Lemma 1 in order of appearance, whereas the article prints the same items with the numbering of its own full document.  The complete native source of the article as distributed in the arXiv e-print for this exact version is bundled unchanged as \texttt{sources/191551/source0.tex}.
\begin{lem} \label{lem:automorphisms_quadric_cone}
Let $L$ be a line on the two-dimensional quadric cone $Q=\VPthree(x_0x_2-x_1^2)\sub \Pthree$. Then there is an automorphism $\alpha \in \Aut(Q)$ such that $\alpha(L)=\VPthree(x_1,x_2)\sub Q$. 
\end{lem}

\begin{proof}
We first construct a subgroup of $\Aut(Q)$ from which we will choose the automorphism $\alpha$. Recall that the singular quadric cone $Q\sub\Pthree$ is uniquely given as the union of lines in $\Pthree$ through the singular point $p:=[0:0:0:1]$ and one point on the planar conic $\Gamma:=\VPtwo(x_0x_2-x_1^2)$ embedded into $Q$. Furthermore, the conic $\Gamma$ is isomorphic to $\Pone$ via the isomorphism 
$$ \Pone \rightarrow \Gamma, \: [u:v]\mapsto [u^2:uv:v^2].$$
This isomorphism induces an inclusion
$$\begin{array}{c c c c c c}
\iota: &\PGL_2(\kk) = \Aut(\Pone) &\overset{\simeq}{\longrightarrow} &\Aut(\Gamma)\sub\GL_3(\kk) &\hookrightarrow &\Aut(Q)\sub \GL_4(\kk)\\
&A=\begin{pmatrix}
a & b\\
c & d
\end{pmatrix}
& \mapsto 
&M_A:=\frac{1}{ad-bc} \begin{pmatrix}
a^2 & 2ab & b^2\\
ac & ad+bc & bd\\
c^2 & 2cd & d^2
\end{pmatrix}
&\mapsto 
&\begin{pmatrix}
& & &0\\
& M_A& &0\\
& & &0\\
0&0&0&1
\end{pmatrix}
\end{array}. $$
This inclusion induces an action of $\PGL_2(\kk)$ on $Q$, which is given by $\PGL_2(\kk)\times Q \rightarrow Q, \: (A,x)\mapsto \iota(A)(x)$. This action fixes the singular point 
$p:=[0:0:0:1]$ of $Q$
and acts transitively on the points of the conic $\Gamma'=\VPthree(x_0x_2-x_1^2,x_3)\simeq \Gamma$. It then acts transitively on the lines in $Q$, which are lines through the singular point $p$ and one point of $\Gamma'$. In particular, we can find an isomorphism $\alpha \in \iota(\PGL_2(\kk))\sub \Aut(Q)$ such that $\alpha(L)=\VPthree(x_1,x_2)$.
\end{proof}

\end{document}
